%! Solving Radical Equations \& Extraneous Solutions — Unit 6, Lesson 6.5 — Homework
\documentclass[10pt]{article}
\usepackage{algebra2-article}
\usepackage{algebra2-boxes}
\newcommand{\TallMath}[1]{$\displaystyle #1\rule[-1.4em]{0pt}{3.2em}$}
\newcommand{\lgrid}[4]{%
  \draw[step=1, linegray!40, very thin] (#1,#3) grid (#2,#4);
  \draw[->, semithick, charcoal] (#1,0)--(#2,0) node[below right, font=\scriptsize]{$x$};
  \draw[->, semithick, charcoal] (0,#3)--(0,#4) node[left, font=\scriptsize]{$y$};}

\begin{document}

\pageheader{Unit 6, Lesson 6.5}{Homework}
\namedateperiod

\vspace{0.10in}

\begin{notesbox}{Practice}
Isolate, raise, solve, \textbf{check}. Every answer needs a check shown, and ``no real solution'' is
a complete answer when it is the right one.

\begin{enumerate}[leftmargin=1.9em, itemsep=11pt]
  \item \textbf{One radical.} Solve each equation and check.
    \begin{center}
    \renewcommand{\arraystretch}{1.8}
    \begin{tabularx}{\linewidth}{@{}X X X@{}}
    (a) \; $\sqrt{x+6}=4$ \newline $x=\blank{1.6cm}$ &
    (b) \; $\sqrt{2x-3}+1=6$ \newline $x=\blank{1.6cm}$ &
    (c) \; $3\sqrt{x}-4=8$ \newline $x=\blank{1.6cm}$ \\
    (d) \; $\sqrt[3]{x+5}=3$ \newline $x=\blank{1.6cm}$ &
    (e) \; $\sqrt[3]{4x}=-2$ \newline $x=\blank{1.6cm}$ &
    (f) \; $\sqrt{x-2}+9=5$ \newline $x=\blank{1.6cm}$ \\
    \end{tabularx}
    \end{center}
    Which one of these six needed no algebra at all, and how did you know? \blank{6.4cm}

  \item \textbf{Variable on both sides.} These are the ones that need sorting. List \emph{every}
        candidate, check each, then give the solution set.
    \begin{center}
    \renewcommand{\arraystretch}{1.7}
    \begin{tabularx}{\linewidth}{@{}l X X X@{}}
    \hline
    \textbf{Equation} & \textbf{Candidates} & \textbf{Extraneous?} & \textbf{Solution set} \\
    \hline
    (a) \; $\sqrt{7-x}=x-1$ & \blank{2.2cm} & \blank{2.2cm} & \blank{2.0cm} \\
    (b) \; $\sqrt{4x+13}=x+4$ & \blank{2.2cm} & \blank{2.2cm} & \blank{2.0cm} \\
    (c) \; $x=\sqrt{x+12}$ & \blank{2.2cm} & \blank{2.2cm} & \blank{2.0cm} \\
    \hline
    \end{tabularx}
    \end{center}
    \vspace{1.30in}
    Item (b) is a warning about a habit. Both of its candidates are \emph{negative numbers}, and
    both of them are genuine solutions. So what is it that actually makes a candidate extraneous?
    \blank{6.0cm}

  \item \textbf{Rational exponents.} Raise both sides to the reciprocal power. Watch for a hidden
        $\pm$.
    \begin{center}
    \renewcommand{\arraystretch}{1.8}
    \begin{tabularx}{\linewidth}{@{}X X X X@{}}
    (a) \; $x^{3/2}=8$ \newline $x=\blank{1.4cm}$ &
    (b) \; $x^{2/3}=25$ \newline $x=\blank{1.8cm}$ &
    (c) \; $(x-1)^{1/3}=4$ \newline $x=\blank{1.4cm}$ &
    (d) \; $x^{4/3}=16$ \newline $x=\blank{1.8cm}$ \\
    \end{tabularx}
    \end{center}
    Two of these have \emph{two} solutions. What do those two have in common? \blank{5.4cm}

  \item \textbf{Read the answers off a graph.} The curve below is $y=\sqrt{x+2}$.
    \begin{center}
    \begin{tikzpicture}[scale=0.42]
      \lgrid{-4}{15}{-2}{5}
      \begin{scope}
        \clip (-4,-2) rectangle (15,5);
        \draw[very thick, forest, domain=-2:15, samples=110, smooth] plot (\x,{sqrt((\x)+2)});
        \draw[very thick, navy] (-4,3)--(15,3);
        \draw[very thick, goldacc] (-4,1)--(15,1);
        \draw[very thick, redacc] (-4,-1)--(15,-1);
      \end{scope}
      \fill[forest] (-2,0) circle (4pt);
      \fill[charcoal] (7,3) circle (3pt);
      \fill[charcoal] (-1,1) circle (3pt);
      \node[navy, font=\scriptsize, right] at (12.4,3.5) {$y=3$};
      \node[goldacc, font=\scriptsize, right] at (12.4,1.5) {$y=1$};
      \node[redacc, font=\scriptsize, right] at (12.4,-1.5) {$y=-1$};
    \end{tikzpicture}
    \end{center}
    \begin{enumerate}[label=(\alph*), leftmargin=1.9em, itemsep=4pt]
      \item $\sqrt{x+2}=3$ has \blank{1.4cm} solution(s): $x=\blank{1.4cm}$
      \item $\sqrt{x+2}=1$ has \blank{1.4cm} solution(s): $x=\blank{1.4cm}$
      \item $\sqrt{x+2}=-1$ has \blank{1.4cm} solution(s), because \blank{4.4cm}
      \item Solve (c) algebraically anyway --- square both sides and see what number comes out.
            $x=\blank{1.4cm}$. \; What is that number called, and what does the graph say about it?
            \blank{5.0cm}
    \end{enumerate}

  \item \textbf{Error analysis.} A classmate turns in this work.
    \begin{center}
    \begin{tcolorbox}[colback=white, colframe=charcoal!45, boxrule=0.7pt, arc=1.5mm,
        width=0.62\linewidth, left=3mm, right=3mm, top=1.5mm, bottom=1.5mm]
    \centering\small
    \renewcommand{\arraystretch}{1.3}
    \begin{tabular}{@{}l@{}}
    $\sqrt{x}+2=5$ \\
    $x+4=25$ \\
    $x=21$ \\
    \end{tabular}
    \end{tcolorbox}
    \end{center}
    \vspace{-4pt}
    (a) Check $x=21$ in the original equation. Does it work? \blank{3.0cm} \\[3pt]
    (b) Which \emph{step} of the method did they skip, and what should line 2 have been?
    \blank{6.0cm} \\[3pt]
    (c) Solve it correctly. $x=\blank{1.6cm}$

  \item \textbf{Explain.} For a \emph{square root} equation the check can reject an answer, but for
        a \emph{cube root} equation the raising step can never produce one to reject. Explain both
        halves of that sentence, using the words \textbf{index} and \textbf{sign}.
        \writelines{4}
\end{enumerate}
\end{notesbox}

\begin{extensionbox}
\textbf{Reading the skid marks.} Crash investigators estimate how fast a car was going from the
length of the skid marks it left. If a car skids $d$ feet on dry asphalt before stopping, its speed
when the brakes locked was about
\[
  s(d)=\sqrt{24d} \ \text{ miles per hour.}
\]
In Lesson 6.0 you were given $d$ and found $s$. An investigator has the opposite problem: they know
the speed they are testing and need the skid length --- which means solving a radical equation.
\begin{enumerate}[label=(\alph*), leftmargin=1.8em, itemsep=5pt]
  \item A driver claims they were doing the $30$ mph speed limit. Solve $30=\sqrt{24d}$ to find how
        long the skid marks would be at that speed.
  \item Now solve for a speed of $60$ mph. Compare the two skid lengths: the speed doubled, but what
        happened to the distance? Which feature of the graph of $s(d)=\sqrt{24d}$ predicted that
        before you did any algebra?
  \item The actual skid marks at the scene are $96$ feet long. How fast was the car going? Was the
        driver telling the truth?
  \item Solving $s=\sqrt{24d}$ for $d$ in general gives $d=\dfrac{s^{2}}{24}$ --- and this is exactly
        the formula you derived in the Lesson 6.0 homework. You obtained it by squaring both sides,
        which is the step that usually manufactures extraneous solutions. Explain why there is no
        risk of one here.
  \item What does $s(0)=0$ say about the situation, and why does the context rule out negative
        values of $d$ even though the algebra would happily accept them?
\end{enumerate}
\writelines{5}
\end{extensionbox}

\begin{spiralbox}
\textbf{Coming next --- Lesson 6.6: Composition of Functions.} Today you undid a radical by raising
both sides to a power --- one operation cancelling another. That idea is about to get a name and a
notation. In 6.6 you will feed one function's output straight into another function,
$(f\circ g)(x)=f(g(x))$, and find out that the order matters enormously. Then in 6.7 composition
becomes the \emph{test}: two functions are inverses exactly when composing them in either order
hands you back the $x$ you started with --- which will finally explain the reflection over $y=x$
you first noticed all the way back in Lesson 6.0.
\end{spiralbox}

\end{document}
